\begin{definitionbox}{Definition (Extended Real Interval)}
\begin{definition}[Extended Real Interval]
\label{def:extended-real-interval}
A set $I\subseteq\overline{\mathbb{R}}$ is an \emph{extended real interval} if
whenever $x,z\in I$ and $x\le y\le z$ in $\overline{\mathbb{R}}$, then
$y\in I$.
\end{definition}
\end{definitionbox}
\begin{remark*}[Standard quantified statement]
\[
\operatorname{ExtInterval}(I)
\Longleftrightarrow
\forall x,z\in I\;\forall y\in\overline{\mathbb{R}}\;
\bigl(x\le y\le z\Rightarrow y\in I\bigr).
\]
\end{remark*}
\begin{remark*}[Predicate reading]
\[
\operatorname{ExtendedRealInterval}(I,\overline{\mathbb{R}},\le)
\]
means that $I$ is convex in the extended real order.
\end{remark*}
\begin{remark*}[Negated quantified statement]
\[
\exists x,z\in I\;\exists y\in\overline{\mathbb{R}}\;
\bigl(x\le y\le z\land y\notin I\bigr).
\]
\end{remark*}

\begin{remark*}[Negation predicate reading]
\[
\exists x,z\in I\;\exists y\in\overline{\mathbb{R}}\;
\bigl(x\le y\le z\land y\notin I\bigr).
\]
\end{remark*}

\begin{remark*}[Failure modes]
\begin{description}
  \item[Exposition.] The named branch below records the displayed negation of the predicate.
  \item[\(\operatorname{ExtendedRealInterval}\) fails.] The predicate fails exactly in the displayed negation case:
  \[
  \exists x,z\in I\;\exists y\in\overline{\mathbb{R}}\;
  \bigl(x\le y\le z\land y\notin I\bigr).
  \]
  In predicate-reading form this is recorded as
  \[
  \exists x,z\in I\;\exists y\in\overline{\mathbb{R}}\;
  \bigl(x\le y\le z\land y\notin I\bigr).
  \]
\end{description}
\end{remark*}
\begin{remark*}[Interpretation]
The interval condition is unchanged, but the ambient ordered set now has two
extra endpoints.
\end{remark*}
\begin{dependencies}
\hyperref[def:extended-real-order]{Order on the Extended Reals};
\hyperref[def:interval]{Interval}.
\end{dependencies}

\begin{theorembox}{Theorem (Real Rays as Extended Intervals)}
\begin{theorem}[Real Rays as Extended Intervals]
\label{thm:real-rays-as-extended-intervals}
For $a\in\mathbb{R}$, the real rays
\[
(a,+\infty),\qquad [a,+\infty),\qquad
(-\infty,a),\qquad (-\infty,a]
\]
are the finite parts of extended intervals with infinite endpoints.
\hyperref[prf:real-rays-as-extended-intervals]{Go to proof.}
\end{theorem}
\end{theorembox}
\begin{remark*}[Standard quantified statement]
\[
[a,+\infty)_{\mathbb{R}}
= [a,+\infty]_{\overline{\mathbb{R}}}\cap\mathbb{R}
\]
and similarly for the other ray types.
\end{remark*}
\begin{remark*}[Interpretation]
What was unbounded in $\mathbb{R}$ becomes endpoint-bounded in
$\overline{\mathbb{R}}$.
\end{remark*}
\begin{dependencies}
\hyperref[def:ray]{Ray};
\hyperref[def:extended-real-interval]{Extended Real Interval}.
\end{dependencies}

\begin{theorembox}{Theorem (Extended Closed Rays Contain Infinite Endpoints)}
\begin{theorem}[Extended Closed Rays Contain Infinite Endpoints]
\label{thm:extended-closed-rays-contain-infinite-endpoints}
For $a\in\mathbb{R}$,
\[
+\infty\in [a,+\infty]_{\overline{\mathbb{R}}}
\quad\text{and}\quad
-\infty\in [-\infty,a]_{\overline{\mathbb{R}}}.
\]
\hyperref[prf:extended-closed-rays-contain-infinite-endpoints]{Go to proof.}
\end{theorem}
\end{theorembox}
\begin{remark*}[Standard quantified statement]
\[
\forall a\in\mathbb{R}\;
\bigl(+\infty\in [a,+\infty]_{\overline{\mathbb{R}}}
\land -\infty\in [-\infty,a]_{\overline{\mathbb{R}}}\bigr).
\]
\end{remark*}
\begin{remark*}[Interpretation]
Closedness in the extended line includes the newly added endpoint.
\end{remark*}
\begin{dependencies}
\hyperref[def:extended-real-order]{Order on the Extended Reals}.
\end{dependencies}

\begin{theorembox}{Theorem (The Extended Real Line Is Closed and Bounded)}
\begin{theorem}[The Extended Real Line Is Closed and Bounded]
\label{thm:extended-real-line-is-closed-bounded}
In $\overline{\mathbb{R}}$,
\[
\overline{\mathbb{R}}=[-\infty,+\infty].
\]
Thus the whole extended real line has both a least element and a greatest
element.
\hyperref[prf:extended-real-line-is-closed-bounded]{Go to proof.}
\end{theorem}
\end{theorembox}
\begin{remark*}[Standard quantified statement]
\[
\forall x\in\overline{\mathbb{R}}\;(-\infty\le x\le+\infty).
\]
\end{remark*}
\begin{remark*}[Interpretation]
The extended line is order-bounded by construction.
\end{remark*}
\begin{dependencies}
\hyperref[thm:comparison-with-infinities]{Comparison with Infinities};
\hyperref[def:extended-real-interval]{Extended Real Interval}.
\end{dependencies}

\begin{theorembox}{Theorem (Extended Interval Order Characterization)}
\begin{theorem}[Extended Interval Order Characterization]
\label{thm:extended-interval-order-characterization}
A subset $I\subseteq\overline{\mathbb{R}}$ is an extended real interval if and
only if for all $x,z\in I$ with $x\le z$,
\[
[x,z]_{\overline{\mathbb{R}}}\subseteq I.
\]
\hyperref[prf:extended-interval-order-characterization]{Go to proof.}
\end{theorem}
\end{theorembox}
\begin{remark*}[Standard quantified statement]
\[
\operatorname{ExtInterval}(I)
\Longleftrightarrow
\forall x,z\in I\;\bigl(x\le z\Rightarrow
[x,z]_{\overline{\mathbb{R}}}\subseteq I\bigr).
\]
\end{remark*}
\begin{remark*}[Predicate reading]
\[
\operatorname{ExtendedRealInterval}(I,\overline{\mathbb{R}},\le)
\Longleftrightarrow
\forall x,z\in I\;(x\le z\Rightarrow [x,z]_{\overline{\mathbb{R}}}\subseteq I).
\]
\end{remark*}
\begin{remark*}[Negated quantified statement]
\[
\exists x,z\in I\;\bigl(x\le z\land [x,z]_{\overline{\mathbb{R}}}\nsubseteq I\bigr).
\]
\end{remark*}

\begin{remark*}[Negation predicate reading]
\[
  \exists x,z\in I\;\bigl(x\le z\land [x,z]_{\overline{\mathbb{R}}}\nsubseteq I\bigr).
\]
\end{remark*}

\begin{remark*}[Failure modes]
\begin{description}
  \item[Exposition.] The named branch below records the displayed negation of the predicate.
  \item[\(\operatorname{ExtendedRealInterval}\) fails.] The predicate fails exactly in the displayed negation case:
  \[
  \exists x,z\in I\;\bigl(x\le z\land [x,z]_{\overline{\mathbb{R}}}\nsubseteq I\bigr).
  \]
  In predicate-reading form this is recorded as
  \[
    \exists x,z\in I\;\bigl(x\le z\land [x,z]_{\overline{\mathbb{R}}}\nsubseteq I\bigr).
  \]
\end{description}
\end{remark*}
\begin{remark*}[Interpretation]
The familiar interval characterization survives after adjoining the two
order endpoints.
\end{remark*}
\begin{dependencies}
\hyperref[def:extended-real-interval]{Extended Real Interval};
\hyperref[thm:interval-characterization]{Characterization of Intervals}.
\end{dependencies}
